\section{The finite CM certificate and the identity tables}
\label{app:cm-certificate}

The arithmetic theorem reduces enumeration to orders of class number at
most two. This appendix specifies the resulting finite certificate.
For an order of discriminant $D$, let $H_D$ be its monic ring class
polynomial. At levels two and three write
\[
 P_\ell(J,C)=J^2-T_\ell(C)J+N_\ell(C).
\]
If $h(D)=2$, irreducibility gives the particularly short test
\[
 P_\ell(J,C)=H_D(J)
 \quad\Longleftrightarrow\quad
 T_\ell(C)=\operatorname{Tr}(j_D),\quad
 N_\ell(C)=\operatorname{Nm}(j_D).
\]
Thus equality of two degree-two characteristic polynomials replaces
expanded resultant calculations. For $h(D)=1$, substitution of its
integer root gives the corresponding one-variable test. Exact
factorization, rather than a bound on the size of $C$, exhausts their
integer solutions.

The conductor verification and defining correspondences are stated in
Section~\ref{sec:enumeration}. Its finite scan gives exactly
\begin{align*}
\mathcal D_1={}&\{-3,-4,-7,-8,-11,-12,-16,-19,-27,-28,-43,-67,-163\},\\
\mathcal D_2={}&\{-15,-20,-24,-32,-35,-36,-40,-48,-51,-52,-60,-64,\\
&-72,-75,-88,-91,-99,-100,-112,-115,-123,-147,-148,\\
&-187,-232,-235,-267,-403,-427\}.
\end{align*}
Here $\mathcal D_i$ denotes the discriminants of orders with class number
$i$. Nonmaximal orders are included.
The executed certificate uses Arb complex intervals to certify the
integer class-polynomial coefficients and independently compares them
with PARI. It checks the integer roots both by exact factorization and
by resultants. These computations certify the finite algebraic step;
the period and modular arguments prove the evaluations.

\input{verified/cm_summary.tex}

\subsection{Standard classes}
The nome and its parameter $N$ use the convention in
\eqref{eq:alpha}. We choose $B>0$. Every row is ordinarily convergent;
the row $C=-64$ at level four is conditional, and all other rows are
absolutely convergent.

\begin{identitytable}\centering\small
\caption{Standard level 1 identities, normalized by $B>0$.}
\label{tab:standard1}
\setlength{\tabcolsep}{4pt}
\renewcommand{\arraystretch}{1.15}
\begin{tabular}{@{}crrrrl@{}}\toprule
nome & $N$ & $A$ & $B$ & $C$ & $\alpha$\\\midrule
+&2&3&28&8000&$5 \sqrt{5}$\\
+&3&1&11&54000&$\frac{5 \sqrt{15}}{6}$\\
+&4&5&63&287496&$\frac{11 \sqrt{33}}{4}$\\
+&7&8&133&16581375&$\frac{85 \sqrt{255}}{54}$\\
-&7&8&63&-3375&$5 \sqrt{15}$\\
-&11&15&154&-32768&$32 \sqrt{2}$\\
-&19&25&342&-884736&$32 \sqrt{6}$\\
-&27&31&506&-12288000&$\frac{160 \sqrt{30}}{9}$\\
-&43&263&5418&-884736000&$\frac{640 \sqrt{15}}{3}$\\
-&67&10177&261702&-147197952000&$1760 \sqrt{330}$\\
-&163&13591409&545140134&-262537412640768000&$426880 \sqrt{10005}$\\
\bottomrule\end{tabular}\end{identitytable}
\begin{identitytable}\centering\small
\caption{Standard level 2 identities, normalized by $B>0$.}
\label{tab:standard2}
\setlength{\tabcolsep}{4pt}
\renewcommand{\arraystretch}{1.15}
\begin{tabular}{@{}crrrrl@{}}\toprule
nome & $N$ & $A$ & $B$ & $C$ & $\alpha$\\\midrule
+&2&1&7&648&$\frac{9}{2}$\\
+&3&1&8&2304&$2 \sqrt{3}$\\
+&5&1&10&20736&$\frac{9 \sqrt{2}}{4}$\\
+&9&3&40&614656&$\frac{49 \sqrt{3}}{9}$\\
+&11&19&280&2509056&$18 \sqrt{11}$\\
+&29&1103&26390&24591257856&$\frac{9801 \sqrt{2}}{4}$\\
-&5&3&20&-1024&$8$\\
-&7&8&65&-3969&$9 \sqrt{7}$\\
-&9&3&28&-12288&$\frac{16 \sqrt{3}}{3}$\\
-&13&23&260&-82944&$72$\\
-&25&41&644&-6635520&$\frac{288 \sqrt{5}}{5}$\\
-&37&1123&21460&-199148544&$3528$\\
\bottomrule\end{tabular}\end{identitytable}

\begin{identitytable}\centering\small
\caption{Standard level 3 identities, normalized by $B>0$.}
\label{tab:standard3}
\setlength{\tabcolsep}{4pt}
\renewcommand{\arraystretch}{1.15}
\begin{tabular}{@{}crrrrl@{}}\toprule
nome & $N$ & $A$ & $B$ & $C$ & $\alpha$\\\midrule
+&2&1&6&216&$3 \sqrt{3}$\\
+&4&2&15&1458&$\frac{27}{4}$\\
+&5&4&33&3375&$\frac{15 \sqrt{3}}{2}$\\
-&9&1&5&-192&$\frac{4 \sqrt{3}}{3}$\\
-&17&7&51&-1728&$12 \sqrt{3}$\\
-&25&1&9&-8640&$\frac{4 \sqrt{15}}{5}$\\
-&41&53&615&-110592&$96 \sqrt{3}$\\
-&49&13&165&-326592&$\frac{108 \sqrt{7}}{7}$\\
-&89&827&14151&-27000000&$1500 \sqrt{3}$\\
\bottomrule\end{tabular}\end{identitytable}
\begin{identitytable}\centering\small
\caption{Standard level 4 identities, normalized by $B>0$.}
\label{tab:standard4}
\setlength{\tabcolsep}{4pt}
\renewcommand{\arraystretch}{1.15}
\begin{tabular}{@{}crrrrl@{}}\toprule
nome & $N$ & $A$ & $B$ & $C$ & $\alpha$\\\midrule
+&3&1&6&256&$4$\\
+&7&5&42&4096&$16$\\
-&2&1&4&-64&$2$\\
-&4&1&6&-512&$2 \sqrt{2}$\\
\bottomrule\end{tabular}\end{identitytable}


\subsection{Linear Euler companions}
The coefficient transport of Theorem~\ref{thm:euler} gives the following
primitive representatives and signed multipliers. Its boundary argument
is necessary: it is the reason that 36 standard classes give 35
ordinary Euler classes.

\begin{identitytable}\centering\footnotesize
\caption{Euler companions of the standard level 1 identities.}
\label{tab:euler1}
\setlength{\tabcolsep}{3pt}
\begin{tabular}{@{}rrrl@{}}\toprule
$a$ & $b$ & $C$ & $\beta$\\\midrule
-33&196&8000&$\frac{625 \sqrt{5}}{14}$\\
7&121&54000&$\frac{625 \sqrt{15}}{66}$\\
97&1323&287496&$\frac{14641 \sqrt{33}}{252}$\\
36872&614061&16581375&$\frac{52200625 \sqrt{255}}{7182}$\\
88&189&-3375&$\frac{625 \sqrt{15}}{63}$\\
159&1078&-32768&$\frac{16384 \sqrt{2}}{77}$\\
77&1026&-884736&$\frac{16384 \sqrt{6}}{171}$\\
7861&128018&-12288000&$\frac{10240000 \sqrt{30}}{2277}$\\
49709&1024002&-884736000&$\frac{327680000 \sqrt{15}}{8127}$\\
6625229&170368002&-147197952000&$\frac{149923840000 \sqrt{330}}{130851}$\\
7575892150829&303862746112002&-262537412640768000&$\frac{64856464530145280000 \sqrt{10005}}{272570067}$\\
\bottomrule\end{tabular}\end{identitytable}
\begin{identitytable}\centering\footnotesize
\caption{Euler companions of the standard level 2 identities.}
\label{tab:euler2}
\setlength{\tabcolsep}{3pt}
\begin{tabular}{@{}rrrl@{}}\toprule
$a$ & $b$ & $C$ & $\beta$\\\midrule
-25&49&648&$\frac{729}{14}$\\
0&1&2304&$\frac{9 \sqrt{3}}{32}$\\
7&80&20736&$\frac{729 \sqrt{2}}{40}$\\
179&2400&614656&$\frac{117649 \sqrt{3}}{360}$\\
83&1225&2509056&$\frac{88209 \sqrt{11}}{1120}$\\
4014919&96059600&24591257856&$\frac{941480149401 \sqrt{2}}{105560}$\\
7&20&-1024&$\frac{32}{5}$\\
776&4225&-3969&$\frac{35721 \sqrt{7}}{65}$\\
25&196&-12288&$\frac{256 \sqrt{3}}{7}$\\
119&1300&-82944&$\frac{23328}{65}$\\
6605&103684&-6635520&$\frac{1492992 \sqrt{5}}{161}$\\
162839&3111700&-199148544&$\frac{2744515872}{5365}$\\
\bottomrule\end{tabular}\end{identitytable}
\begin{identitytable}\centering\footnotesize
\caption{Euler companions of the standard level 3 identities.}
\label{tab:euler3}
\setlength{\tabcolsep}{3pt}
\begin{tabular}{@{}rrrl@{}}\toprule
$a$ & $b$ & $C$ & $\beta$\\\midrule
-5&6&216&$6 \sqrt{3}$\\
4&75&1458&$\frac{729}{20}$\\
32&363&3375&$\frac{1875 \sqrt{3}}{22}$\\
14&25&-192&$\frac{64 \sqrt{3}}{15}$\\
10&51&-1728&$\frac{192 \sqrt{3}}{17}$\\
10&81&-8640&$\frac{64 \sqrt{15}}{9}$\\
268&3075&-110592&$\frac{98304 \sqrt{3}}{205}$\\
718&9075&-326592&$\frac{46656 \sqrt{7}}{55}$\\
43834&750003&-27000000&$\frac{375000000 \sqrt{3}}{4717}$\\
\bottomrule\end{tabular}\end{identitytable}
\begin{identitytable}\centering\footnotesize
\caption{Euler companions of the standard level 4 identities.}
\label{tab:euler4}
\setlength{\tabcolsep}{3pt}
\begin{tabular}{@{}rrrl@{}}\toprule
$a$ & $b$ & $C$ & $\beta$\\\midrule
-1&6&256&$\frac{16}{3}$\\
13&126&4096&$\frac{1024}{21}$\\
5&18&-512&$\frac{16 \sqrt{2}}{3}$\\
\bottomrule\end{tabular}\end{identitytable}

